%%%%%%%%%%%%%%%%%%%%%%%%%%%%%%%%%%%%% C A B E C E R A %%%%%%%%%%%%%%%%%%%%%%%%%%%%%%%%%%
% Definimos el estilo del documento
\documentclass[12pt, letterpaper]{report}

% Utilizamos un paquete para gestionar acentos y e\~nes
\usepackage[latin1]{inputenc}
\usepackage[T1]{fontenc}
\usepackage{xcolor}
\usepackage{pifont}
% Utilizamos el paquete para gestionar imagenes jpg
\usepackage{graphicx}
\graphicspath{ {images/} }
% Letra capiral
\usepackage{lettrine}
\usepackage{enumitem}
\usepackage{lmodern}

% Definimos la zona de la pagina ocupada por el texto
\oddsidemargin -1.0cm
\headsep -1.0cm
\textwidth=18.5cm
\textheight=23cm

%SALTO DE PAGINA%%%
%\newpage

%Empieza el documento %%%%%%%%%%%%%%%%%%%% P R I N C I P I O %%%%%%%%%%%%%%%%%%%%%%%%%%%%%%
\begin{document}

% Definimos TÍTULO
\title{
    \Huge \bfseries Sacramento del Bautismo \\
    \large Celebraci\'on de los Sacramentos de la Iniciaci\'on Cristiana \\
    \vspace{2em}
    \includegraphics{cruz_cristo}
}
\author{
    \Large \em \bfseries Martina Caridad Guti\'errez Martinez \\
}
\date{\bfseries 30 de Enero, 2020}

% Generamos t\'itulo
\maketitle

\begin{center}
\Large {\bfseries \textcolor{red}{30 de Enero}}
\end{center}

\begin{center}
\Huge {\bfseries Santa Martina, Virgen y Martir} \\
\Large {\bfseries \em \textcolor{red}{Memoria, Com\'un de una Virgen Martir}}
\end{center}

\Large {Santa Martina, noble virgen romana, luego de padecer atroces suplicios, fue decapitada en Roma en el a\~no 228.}

\begin{center}
\Huge {\bfseries Ritos Iniciales}
\end{center}

\Large {\bfseries \textcolor{red}{ANT\'IFONA DE ENTRADA} } \newline

\Large {Esta es la virgen dichosa que, neg\'andose a s\'i misma y abrazando su cruz, sigui\'o al Se\~nor, esposo de las v\'irgenes y pr\'incipe de los m\'artires.} \newline

% \Large {\bfseries \textcolor{red}{Se dice} Gloria.} \newline

\Large {\bfseries \textcolor{red}{ORACI\'ON COLECTA}}

\lettrine[lines=1]{\bfseries \textcolor{red}{D}}{}\Large {ios,
que nos alegras hoy \\
con la conmemoraci\'on anual de Santa {\bfseries \textcolor{red}{Martina}}, \\
concede la ayuda de sus m\'eritos \\
a cuantos hemos sido iluminados \\
con el ejemplo de su virginidad y fortaleza. \\
Por nuestro Se\~nor Jesucristo, tu Hijo, \\
que vive y reina contigo en la unidad del Esp\'iritu Santo \\
y es Dios por los siglos de los siglos.}

%SALTO DE PAGINA%%%
\newpage

\begin{center}
\Huge {\bfseries Lit\'urgia de la Palabra}
\end{center}

\Large {\bfseries Lectura del libro del Cantar de los cantares \hspace{1cm} \textcolor{red}{8, 6-7}} \newline

\Large Gr\'abame como un sello en tu brazo, como un sello en tu coraz\'on, porque es fuerte el amor como la muerte, es cruel la pasi\'on como el abismo; es centella de fuego, llamarada divina: las aguas torrenciales no podr\'an apagar el amor, ni anegarlo los r\'ios. \\
Si alguien quisiera comprar el amor con todas las riquezas de su casa, se har\'ia despreciable. \newline

{\bfseries Palabra de Dios.} \newline

%SALTO DE PAGINA%%%
% \newpage

\Large {\bfseries \textcolor{red}{Salmo responsorial \hspace{1cm} 44, 11-12. 14-15. 16-17 (R/.: Mt 25, 6b)}} \newline

\Large {\bfseries \textcolor{red}{R/.}} \hspace{1cm} Llega el Esposo; salid a recibir a Cristo, el Se\~nor. \newline

{\bfseries \textcolor{red}{V/.}} \hspace{1cm} Escucha, hija, mira: inclina el o\'ido, \\
. \hspace{2.5cm} olvida tu pueblo y la casa paterna; \\
. \hspace{2.5cm} prendado est\'a el rey de tu belleza: \\
. \hspace{2.5cm} p\'ostrate ante \'el, que \'el es tu Se\~nor. \hspace{1cm} {\bfseries \textcolor{red}{R/.}} \newline

{\bfseries \textcolor{red}{V/.}} \hspace{1cm} Ya entra la princesa, bell\'isima, \\
. \hspace{2.5cm} vestida de perlas y brocado; \\
. \hspace{2.5cm} la llevan ante el rey, con s\'equito de v\'irgenes \\
. \hspace{2.5cm} la siguen sus compa\~neras. \hspace{1cm} {\bfseries \textcolor{red}{R/.}} \newline

%SALTO DE PAGINA%%%
\newpage

{\bfseries \textcolor{red}{V/.}} \hspace{1cm} Las traen entre alegr\'ia y algazara, \\
. \hspace{2.5cm} van entrando en el palacio real. \\
. \hspace{2.5cm} A cambio de tus padres, tendr\'as hijos, \\
. \hspace{2.5cm} que nombrar\'as pr\'incipes por toda la tierra. \hspace{1cm} {\bfseries \textcolor{red}{R/.}}

\begin{center}
\Large {\bfseries \textcolor{red}{Aleluya \hspace{1cm} Jn 14, 23}} \\
El que me ama guardar\'a mi palabra \\
--dice el Se\~nor--, \\
y mi Padre lo amar\'a, y vendremos a \'el.
\end{center}

%SALTO DE PAGINA%%%
% \newpage

\Huge \textcolor{red}{\ding{64}} \Large {\bfseries Lectura del santo evangelio seg\'un san Lucas \hspace{1cm} \textcolor{red}{10, 38-42}} \newline

\Large En aquel tiempo, entr\'o Jes\'us en una aldea, y una mujer llamada Marta lo recibi\'o en su casa. \'Esta ten\'ia una hermana llamada Mar\'ia, que, sentada a los pies del Se\~nor, escuchaba su palabra. Y Marta se multiplicaba para dar abasto con el servicio; hasta que se par\'o y dijo: \\
--<<Se\~nor, ?`No te importa que mi hermana me haya dejado sola con el servicio? Dile que me eche una mano>>--. \\
Pero el Se\~nor le contest\'o: \\
--<<Marta, Marta, andas inquieta y nerviosa con tantas cosas; s\'olo una es necesaria. Mar\'ia ha escogido la parte mejor, y no se la quitar\'an>>--. \newline

{\bfseries Palabra del Se\~nor.} \newline

\Large {\bfseries \textcolor{red}{HOMILIA}} \newline

%SALTO DE PAGINA%%%
\newpage

\begin{center}
\Huge {\bfseries Lit\'urgia Bautismal}
\end{center}

\vspace{1cm}

\large {\bfseries \textcolor{red}{El sacerdote, con los ministros, se dirige a la fuente bautismal.}} \newline

\Large {\bfseries \textcolor{red}{LETANIAS DE LOS SANTOS}} \newline

\large {\bfseries \textcolor{red}{Dos cantores entonan las letan\'ias a las que todos responden estando en pie. En las letan\'ias se pueden a\~nadir algunos nombres de santos, especialmente el del titular de la iglesia, el de los patronos del lugar y el de los que van a ser bautizados.}} \newline

\begin{tabular}{ll}
Se\~nor, ten piedad.& Se\~nor, ten piedad.\\
Cristo, ten piedad.& Cristo, ten piedad. \\
Se\~nor, ten piedad.& Se\~nor, ten piedad.
\end{tabular} \newline \\

\begin{tabular}{ll}
Santa Mar\'ia, Madre de Dios.& Ruega por nosotros. \\
<<<<<<< HEAD
=======
San Miguel.& Ruega por nosotros. \\
>>>>>>> 7d2a9117e1986b96cf06e7721201920f2591461a
Santos \'Angeles de Dios.& Rogad por nosotros. \\
San Jos\'e.& Ruega por nosotros. \\
San Juan Bautista.& Ruega por nosotros. \\
Santos Pedro y Pablo.& Rogad por nosotros. \\
San Juan.& Ruega por nosotros. \\
Santa Mar\'ia Magdalena.& Ruega por nosotros. \\
Santa Martina.& Ruega por nosotros. \\
San Esteban.& Ruega por nosotros. \\
Santos Gregorio y Le\'on.& Rogad por nosotros. \\
San Agust\'in.& Ruega por nosotros. \\
San Basilio.& Ruega por nosotros. \\
San Juan Mar\'ia Vianney.& Ruega por nosotros. \\
Santos Francisco y Domingo.& Ruegad por nosotros. \\
Santos Isidro Labrado y Mar\'ia de la Cabeza.& Ruegad por nosotros. \\
Santos Luis y Celia Martan.& Ruegad por nosotros. \\
Santa Gema.& Ruega por nosotros. \\
Santo Domingo Sabio.& Ruega por nosotros. \\
Santos y Santas de Dios.& Ruegad por nosotros.
\end{tabular} \newline \\

\begin{tabular}{ll}
Mu\'estrate propicio.& L\'ibranos, Se\~nor. \\
De todo mal.& L\'ibranos, Se\~nor. \\
De todo pecado.& L\'ibranos, Se\~nor. \\
De la muerte eterna.& L\'ibranos, Se\~nor. \\
Por tu encarnaci\'on.& L\'ibranos, Se\~nor. \\
Por tu muerte y resurrecci\'on.& L\'ibranos, Se\~nor. \\
Por el env\'io del Esp\'iritu Santo.& L\'ibranos, Se\~nor.
\end{tabular} \newline \\

\begin{tabular}{ll}
Nosotros, que somos pecadores.& Te rogamos, \'oyenos. \\
Para que regeneres a Martina Caridad \\
con la gracia del bautismo.& Te rogamos, \'oyenos. \\
Para que regeneres a estos elegidos \\
con la gracia del bautismo.& Te rogamos, \'oyenos. \\
Jes\'us, Hijo de Dios vivo.& Te rogamos, \'oyenos.
\end{tabular} \newline \\

\begin{tabular}{ll}
Cristo, \'oyenos.& Cristo, \'oyenos. \\
Cristo, esc\'uchanos.& Cristo, esc\'uchanos. \\
\end{tabular} \newline \\

\Large {\bfseries \textcolor{red}{BENDICI\'ON DEL AGUA BAUTISMAL}} \newline

\large {\bfseries \textcolor{red}{Habi\'endose acercado a la pila bautismal, el Sacerdote recordar\'a brevemente a los presentes la admirable providencia de Dios, que ha querido santificar el alma y el cuerpo del hombre por medio del agua. Lo puede hacer con estas o parecidas palabras:}} \newline

\lettrine[lines=1]{\bfseries \textcolor{red}{S}}{}\Large aben, hermanos, \\ 
que Dios concede a los creyentes \\
la abundancia de su vida por el sacramento del agua. \\
Levantemos hacia \'El nuestro coraz\'on, \\
para que por el agua de esta fuente \\
derrame su gracia sobre {\bfseries \textcolor{red}{Martina Caridad}}. \newline

\large {\bfseries \textcolor{red}{El Sacerdote, vuelto hacia la fuente, dice la siguiente bendici\'on:}} \newline

\lettrine[lines=1]{\bfseries \textcolor{red}{O}}{}\Large h Dios, \\
que realizas en tus sacramentos \\
obras admirables con tu poder invisible, \\
y de diversos modos \\
te has servido del agua, tu creatura, \\
para significar la gracia del bautismo. \newline

\Large Mira, ahora, a tu Iglesia en oraci\'on \\
y abre para ella la fuente de Bautismo: \\
que esta agua reciba, por el Esp\'iritu Santo, \\
la gracia de tu Unig\'enito, \\
para que el hombre, \\
creado a tu imagen y limpio en el Bautismo, \\
muera al hombre viejo y renazca, como ni\~no, \\
a nueva vida por el agua y el Esp\'iritu. \newline

\large {\bfseries \textcolor{red}{El Sacerdote toca el agua con la mano derecha y prosigue:}}

\lettrine[lines=1]{\bfseries \textcolor{red}{T}}{}\Large e pedimos, Se\~nor, \\
que el poder del Esp\'iritu Santo, \\
por tu Hijo, \\
descienda sobre el agua de esta fuente \\
para que los sepultados con Cristo en su muerte, \\
por el Bautismo, resuciten con \'El a la vida. \\
Por Jesucristo nuestro Se\~nor. \newline

\Large {\bfseries \textcolor{red}{RENUNCIAS}} \newline

\Large Queridos padres y padrinos: \newline

\Large En el sacramento del Bautismo, esta ni\~na que han presentado a la Iglesia va a recibir, por el agua y el Esp\'iritu Santo, una nueva vida que brota del amor de Dios.\\
Vosotros, por vuestra parte, deben esforzarse en educarlo en la fe de tal manera que esta vida divina quede preservada del pecado y crezca en \'el de d\'ia en d\'ia. \\ \\
As\'i, pues, si est\'an dispuestos a aceptar esta obligaci\'on, recordando vuestro propio bautismo, renuncien al pecado y confiesen su fe en Cristo Jes\'us, que es la fe de la Iglesia, en la que va a ser bautizado su hijo. \newline

\noindent
\Large {\bfseries \textcolor{red}{V/.}} \hspace{0.5cm} ?`Renuncian a Satan\'as, esto es:
\begin{itemize}[leftmargin=2.0cm,labelsep=0.5cm]
\item al pecado, como negaci\'on de Dios;
\item al mal, como signo del pecado en el mundo;
\item al error, como ofuscaci\'on de la verdad;
\item a la violencia, como contraria a la caridad;
\item al ego\'ismo, como falta de testimonio del amor?
\end{itemize}

\noindent
\Large {\bfseries \textcolor{red}{R/.}} \hspace{0.5cm} S\'i renuncio. \newline

%SALTO DE PAGINA%%%
\newpage

\noindent
\Large {\bfseries \textcolor{red}{V/.}} \hspace{0.5cm} ?`Renuncian a sus obras, que son:
\begin{itemize}[leftmargin=2.0cm,labelsep=0.5cm]
\item sus envidias y odios;
\item sus perezas e indiferencias;
\item sus cobard\'ias y complejos;
\item sus tristezas y desconfianzas;
\item sus materialismos y sensualidades;
\item sus injusticias y favoritismos;
\item sus faltas de fe, de esperanza y de caridad?
\end{itemize}

\Large {\bfseries \textcolor{red}{R/.}} \hspace{0.5cm} S\'i renuncio. \newline

\noindent
\Large {\bfseries \textcolor{red}{V/.}} \hspace{0.5cm} ?`Renuncian a todas sus seducciones, como pueden ser:
\begin{itemize}[leftmargin=2.0cm,labelsep=0.5cm]
\item el creerse superiores;
\item el estar muy seguros de vosotros mismos;
\item el creer que ya est\'ais convertidos del todo?
\end{itemize}
\Large {\bfseries \textcolor{red}{R/.}} \hspace{0.5cm} S\'i renuncio. \newline

%SALTO DE PAGINA%%%
\newpage

\noindent
\Large {\bfseries \textcolor{red}{V/.}} \hspace{0.5cm} ?`Renuncian a los criterios y comportamientos materialistas que consideran:
\begin{itemize}[leftmargin=2.0cm,labelsep=0.5cm]
\item el dinero como la aspiraci\'on suprema de la vida; el placer ante todo;
\item el negocio como valor absoluto;
\item el propio bien por encima del bien com\'un?
\end{itemize}
\Large {\bfseries \textcolor{red}{R/.}} \hspace{0.5cm} S\'i renuncio. \newline

\Large {\bfseries \textcolor{red}{PROFESI\'ON DE FE}} \newline

\large {\bfseries \textcolor{red}{Seguidamente el Sacerdote pide esta triple profesi\'on de fe a los padres y padrinos:}} \newline

\noindent
\Large {\bfseries \textcolor{red}{V/.}} \hspace{0.5cm} ?`Creen en Dios, Padre todopoderoso, Creador del cielo y de la tierra? \\
\Large {\bfseries \textcolor{red}{R/.}} \hspace{0.5cm} S\'i creo. \newline

\noindent
\Large {\bfseries \textcolor{red}{V/.}} \hspace{0.5cm} ?`Creen en Jesucristo, su \'unico Hijo, nuestro Se\~nor, que naci\'o de Santa Mar\'ia Virgen, muri\'o, fue sepultado, resucit\'o de entre los muertos y est\'a sentado a la derecha del Padre? \\
\Large {\bfseries \textcolor{red}{R/.}} \hspace{0.5cm} S\'i creo. \newline

\noindent
\Large {\bfseries \textcolor{red}{V/.}} \hspace{0.5cm} ?`Creen en el Esp\'iritu Santo, en la Santa Iglesia Cat\'olica, en la comuni\'on de los Santos, en el perd\'on de los pecados, en la resurrecci\'on de los muertos y en la vida eterna? \\
\Large {\bfseries \textcolor{red}{R/.}} \hspace{0.5cm} S\'i creo. \newline

%SALTO DE PAGINA%%%
\newpage

\Large {\bfseries \textcolor{red}{BAUTISMO}} \newline

\large {\bfseries \textcolor{red}{El Sacerdote invita a la familia para que se acerque a la fuente, y pregunta a los padres y padrinos:}} \newline

\noindent
\Large {\bfseries \textcolor{red}{V/.}} \hspace{0.5cm} ?`Quieren, por tanto, que su hija sea bautizada en la fe de la Iglesia, que todos juntos acabamos de profesar? \\
\Large {\bfseries \textcolor{red}{R/.}} \hspace{0.5cm} S\'i quiero. \newline

\large {\bfseries \textcolor{red}{Bautismo por inmersi\'on.}} \newline

\Large {\bfseries \textcolor{red}{UNCI\'ON CON EL SANTO CRISMA}}

\lettrine[lines=1]{\bfseries \textcolor{red}{D}}{}\Large ios todopoderoso, \\
Padre de nuestro Se\~nor Jesucristo, \\
que te ha liberado del pecado \\
y dado nueva vida por el agua y el Esp\'iritu Santo, \\
te consagre con el crisma de la salvaci\'on \\
para que entres a formar parte de su pueblo \\
y seas para siempre miembro de Cristo, \\
sacerdote, profeta y rey. \newline

\noindent
\Large {\bfseries \textcolor{red}{R/.}} \hspace{0.5cm} Am\'en. \newline

\large {\bfseries \textcolor{red}{Seguidamente, en silencio, el Sacerdote unge la coronilla de la ni\~no(a) con el santo crisma.}} \newline

%SALTO DE PAGINA%%%
\newpage

\Large {\bfseries \textcolor{red}{IMPOSICI\'ON DE LA VESTIDURA BLANCA}}

\lettrine[lines=1]{\bfseries \textcolor{red}{M}}{}\Large artina Caridad, \\ 
eres ya nueva creatura \\
y has sido revestido de Cristo. \\
Esta vestidura blanca sea signo de tu dignidad de cristiano. \\
Ayudado por la palabra y el ejemplo de los tuyos, \\
cons\'ervala sin mancha hasta la vida eterna. \newline

\noindent
\Large {\bfseries \textcolor{red}{R/.}} \hspace{0.5cm} Am\'en. \newline

\Large {\bfseries \textcolor{red}{ENTREGA DEL CIRIO}} \newline

\large {\bfseries \textcolor{red}{El sacerdote muestra el cirio pascual y dice:}} \newline

\noindent
\Large {\bfseries \textcolor{red}{V/.}} \hspace{0.5cm} Recibe la luz de Cristo. \newline

\large {\bfseries \textcolor{red}{El padrino, enciende la vela del ni\~no(a) en el cirio pascual.}}

\lettrine[lines=1]{\bfseries \textcolor{red}{A}}{} \Large vosotros, padres y padrinos, \\
se les conf\'ia acrecentar esta luz.\\
Que vuestra hija, iluminada por Cristo, \\
camine siempre como hija de la luz. \\
Y perseverando en la fe, \\
pueda salir con todos los Santos \\
al encuentro del Se\~nor. \newline

%SALTO DE PAGINA
\newpage

\Large {\bfseries \textcolor{red}{ \'EFFETA}} \newline

\large {\bfseries \textcolor{red}{El sacerdote , tocando con el dedo pulgar los o\'idos y la boca del ni\~no(a), dice:}}

\lettrine[lines=1]{\bfseries \textcolor{red}{E}}{}\Large l Se\~nor Jes\'us, \\
que hizo o\'ir a los sordos y hablar a los mudos, \\
te conceda, a su tiempo, \\
escuchar su Palabra y proclamar la fe, \\
para alabanza y gloria de Dios Padre. \newline

\noindent
\Large {\bfseries \textcolor{red}{R/.}} \hspace{0.5cm} Am\'en. \newline

%SALTO DE PAGINA%%%
\newpage

\Large {\bfseries \textcolor{red}{ORACI\'ON DE LOS FIELES}} \newline

\large {\bfseries \textcolor{red}{Intenciones:}} \newline

\large {\bfseries \textcolor{red}{a) por las necesidades de la Iglesia;}} \newline

\Large Por la santa Iglesia, por los sacerdotes,religiosos,diaconos y seminaristas para que Dios les conceda cumplir fielmente la mision encomendada por Jesucristo. \hspace{1cm} \textcolor{red}{R/.}\newline

\large {\bfseries \textcolor{red}{b) por los gobernantes y por la salvaci\'on del mundo entero;}} \newline

\Large Por los gobernantes del mundo entero especialmente el de nuestro pais, para que Dios en su infinita misericordia les ilumine a tomar acciones nobles que ayuden a los mas necesitados y al bien comun. \hspace{1cm} \textcolor{red}{R/.}\newline

\large {\bfseries \textcolor{red}{c) por aquellos que se encuentran en necesidades particulares;}} \newline

\Large Por todos que sufren a causa de las enfermedades,guerras y hambruna para que el Se\~nor sea su auxilio y su ayuda. \hspace{1cm} \textcolor{red}{R/.}\newline

\large {\bfseries \textcolor{red}{d) por la comunidad local;}} \newline

\Large Por todos los aqui presentes para que renovados por este sacramento que juntos vamos a celebrar vivan con la alegria de anunciar con sus vidas el reino de Dios. \hspace{1cm} \textcolor{red}{R/.}\newline

\large {\bfseries \textcolor{red}{e) por ni\~no(a) a ser bautizado;}} \newline

\Large Por la nueva bautizada  \textcolor{red}{Martina Caridad}, para que el, Esp\'iritu. Santo la llene con su gracia y la haga vivir plenamente en la fe, la esperanza y la caridad. \hspace{1cm} \textcolor{red}{R/.}

%SALTO DE PAGINA%%%
\newpage

\Large {\bfseries \textcolor{red}{RITO DE LA PAZ}} \newline

\noindent
\Large {\bfseries \textcolor{red}{V/.}} \hspace{0.5cm} La paz del Se\~nor est\'e siempre con ustedes. \\
\Large {\bfseries \textcolor{red}{R/.}} \hspace{0.5cm} Y con tu esp\'iritu. \newline

\noindent
\Large {\bfseries \textcolor{red}{V/.}} \hspace{0.5cm} Dense fraternalmente la paz. \newline

%SALTO DE PAGINA%%%
\newpage

\begin{center}
\Huge {\bfseries Lit\'ugia Eucar\'istica}
\end{center}

\noindent
\Large {\bfseries \textcolor{red}{V/.}} \hspace{0.5cm} Oren, hermanos, \\
. \hspace{1.5cm} para que este sacrificio, m\'io y de ustedes, \\
. \hspace{1.5cm} sea agradable a Dios, Padre todopoderoso. \newline

\noindent
\Large {\bfseries \textcolor{red}{R/.}} \hspace{0.5cm} El Se\~nor reciba de tus manos este sacrificio \\
. \hspace{1.5cm} para alabanza y gloria de su nombre, \\
. \hspace{1.5cm} para nuestro bien \\
. \hspace{1.5cm} y el de toda su santa Iglesia. \newline

\Large {\bfseries \textcolor{red}{ORACI\'ON SOBRE LAS OFRENDAS}}

\lettrine[lines=1]{\bfseries \textcolor{red}{S}}{}\Large e\~nor, que los dones que te presentamos \\
en la fiesta de Santa {\bfseries \textcolor{red}{Martina}} \\
sean tan agradables a tu bondad \\
como lo fue para ti el combate de su martirio. \\
Por Jesucristo, nuestro Se\~nor. \newline

\Large {\bfseries \textcolor{red}{ANT\'IFONA DE LA COMUNI\'ON \hspace{1cm} Ap 7, 17}} \newline

\Large El Cordero que est\'a delante del trono \\
los conducir\'a hacia fuentes de aguas vivas.

%% SALTO DE PAGINA
\newpage

\begin{center}
\Large PLEGARIA EUCARISTICA
\end{center}

\Large {\bfseries \textcolor{red}{PREFACIO}} \newline

\noindent
\Large {\bfseries \textcolor{red}{V/.}} \hspace{0.5cm} El Se\~nor est\'e con vosotros. \\
\Large {\bfseries \textcolor{red}{R/.}} \hspace{0.5cm} Y con tu esp\'iritu. \newline

\noindent
\Large {\bfseries \textcolor{red}{V/.}} \hspace{0.5cm} Levantemos el coraz\'on. \\
\Large {\bfseries \textcolor{red}{R/.}} \hspace{0.5cm} Lo tenemos levantado hacia el Se\~nor. \newline

\noindent
\Large {\bfseries \textcolor{red}{V/.}} \hspace{0.5cm} Demos gracias al Se\~nor, nuestro Dios. \\
\Large {\bfseries \textcolor{red}{R/.}} \hspace{0.5cm} Es justo y necesario.

\lettrine[lines=1]{\bfseries \textcolor{red}{E}}{}\Large n verdad es justo y necesario, \\
es nuestro deber y salvaci\'on \\
darte gracias, siempre y en todo lugar, \\
a ti Padre Santo, por Jesucristo, tu Hijo amado. \newline

Por \'el, que es tu parabra, \\
hiciste todas las cosas; \\
t\'u nos lo enviaste \\
hecho hombre por obra del Esp\'iritu Santo \\
para que nacido de Mar\'ia, la Virgen, \\
fuera nuestro Salvador y Redentor. \newline

\'El, en cumplimiento de tu voluntad, \\
para destruir la muerte y manifestar la resurrecci\'on, \\
extendi\'o sus brazos en la cruz, \\
y as\'i adquiri\'o para ti un pueblo santo. \newline

Muriendo, destruto nuesta muerte, \\
resucitando, restauro nuesta vida. \newline

Por eso, \\
con los \'angeles y los santos, \\
cantamos tu gloria diciendo: \newline

Santo, Santo, Santo...

%SALTO DE PAGINA%%%
\newpage

\Large {\bfseries \textcolor{red}{CONSAGRACI\'ON}} \newline

\large{\bfseries \textcolor{red}{El sacerdote, con las manos extendidas, dice:}}

\lettrine[lines=1]{\bfseries \textcolor{red}{S}}{}\Large anto eres en verdad, Se\~nor,\\
fuente de toda santidad;\newline

\large{\bfseries \textcolor{red}{Junta las manos y, manteni\'endolas extendidas sobre las ofrendas, dice:}}\newline

\Large por eso te pedimos que santifiques estos dones\\
con la efusi\'on de tu Esp\'iritu,\newline

\large{\bfseries \textcolor{red}{Junta las manos y traza el signo de la cruz sobre el pan y el c\'aliz conjuntamente, diciendo:}}\newline

\Large de manera que se conviertan para nosotros\\
en el Cuerpo y \Huge{\textcolor{red}{\ding{64}}} \Large Sangre\\
de Jesucristo, nuestro Se\~nor.\newline

\large{\bfseries \textcolor{red}{Junta las manos.}}

\lettrine[lines=1]{\bfseries \textcolor{red}{E}}{}\Large l cual,\\
cuando iba a ser entregado a su Pasi\'on,\\
voluntariamente aceptada,\newline

\large{\bfseries \textcolor{red}{Toma el pan y, sosteni\'endolo un poco elevado sobre el altar, prosigue:}}\newline

\Large tom\'o pan, d\'andote gracias, lo parti\'o\\
y lo dio a sus disc\'ipulos, diciendo:\newline

\large{\bfseries \textcolor{red}{Se inclina un poco.}} \newline

\LARGE{ \bfseries{ Tomen y coman todos de \'el,\\
porque esto es mi Cuerpo,\\
que ser\'a entregado por ustedes.}}\newline

\large{\bfseries \textcolor{red}{Muestra el pan consagrado al pueblo, lo deposita luego sobre la patena y lo adora haciendo genuflexi\'on. Despu\'es prosigue:}}\newline

\Large Del mismo modo, acabada la cena,\newline

\large{\bfseries \textcolor{red}{Toma el c\'aliz y, sosteni\'endolo un poco elevado sobre el altar, prosigue:}}\newline

\Large tom\'o el c\'aliz,\\
y, d\'andote gracias de nuevo,\\
lo pas\'o a sus disc\'ipulos, diciendo:\newline

\large{\bfseries \textcolor{red}{Se inclina un poco.}}\newline

\LARGE{ \bfseries{ Tomen y beban todos de \'el,\\
porque este es el c\'aliz de mi Sangre,\\
Sangre de la alianza nueva y eterna,\\
que ser\'a derramada\\
por ustedes y por muchos\\
para el perdon de los pecados.\\
Hagan esto en conmemoraci\'on m\'ia.}}\newline

\large{\bfseries \textcolor{red}{Muestra el c\'aliz al pueblo, lo deposita luego sobre el corporal y lo adora haciendo genuflexi\'on.}} \newline \\

\noindent
\Large {\bfseries \textcolor{red}{V/.}} \hspace{0.5cm} \'Este es el Sacramento de nuestra fe.\\
\Large {\bfseries \textcolor{red}{R/.}} \hspace{0.5cm} Anunciamos tu muerte, Se\~nor,\\
. \hspace{1.5cm} proclamamos tu resurrecci\'on.\\
. \hspace{1.5cm} Ven, Se\~nor Jes\'us!\newline

\large{\bfseries \textcolor{red}{Despu\'es el sacerdote, con las manos extendidas, dice:}} \newline

%SALTO DE PAGINA%%%
% \newpage

\Large {\bfseries \textcolor{red}{ANAMNESIS}}

\lettrine[lines=1]{\bfseries \textcolor{red}{A}}{}\Large s\'i, pues, Padre,\\
al celebrar ahora el memorial\\
de la muerte y resurrecci\'on de tu Hijo,\\
te ofrecemos el pan de vida \\
y el c\'aliz de salvaci\'on,\\
y te damos gracias\\
porque nos haces dignos de servirte en tu presencia.

\lettrine[lines=1]{\bfseries \textcolor{red}{T}}{}\Large e pedimos humildemente\\
que el Esp\'iritu Santo congregue en la unidad\\
a cuantos participamos\\
del Cuerpo y Sangre de Cristo.\newline

Acu\'erdate, Se\~nor,\\
de tu Iglesia extendida por toda la tierra; \\
y con el Papa {\bfseries \textcolor{red}{N.}},\\
con nuestro Obispo {\bfseries \textcolor{red}{N.}},\\
y todos los pastores que cuidan de tu pueblo,\\
llevala a su perfecci\'on por la caridad.\newline

Acu\'erdate tambi\'en de {\bfseries \textcolor{red}{N.}},\\
que hoy has hecho renacer \\
del agua y del Esp\'iritu Santo, \\
libr\'andolo del pecado; \\
t\'u que los has incorporado, \\
como miembro vivo, al cuerpo de Cristo, \\
inscribe tambi\'en su nombre en el libro de la vida.

\lettrine[lines=1]{\bfseries \textcolor{red}{A}}{}\Large cu\'erdate tambi\'en de nuestros hermanos\\
que se durmieron en la esperanza\\
de la resurrecc\'ion,\\
y de todos los que han muerto en tu misericordia;\\
admitelos a contemplar la luz de tu rostro.\newline

Ten misericordia de todos nosotros,\\
y as\'i, con Mar\'a, la Virgen Madre de Dios,\\
San Jos\'e su Casto Esposo, los ap\'ostoles\\
y cuantos vivieron en tu amistad\\
a trav\'es de los tiempos,\\
merezcamos, por tu Hijo Jesucristo,\\
compartir la vida eterna\\
y cantar tus alabanzas. \newline

%SALTO DE PAGINA%%%
\newpage


\large{\bfseries \textcolor{red}{Junta las manos. Toma la patena con el pan consagrado y el c\'aliz y, sosteni\'endolos elevados, dice:}}

\lettrine[lines=1]{\bfseries \textcolor{red}{P}}{}\Large or Cristo, con \'el y en \'el,\\
a ti, Dios Padre omnipotente,\\
en la unidad del Esp\'iritu Santo,\\
todo honor y toda gloria\\
por los siglos de los siglos.\newline

\Large \hspace{-0.9cm} {\bfseries \textcolor{red}{R/.}} \hspace{0.5cm} Am\'en.

%SALTO DE PAGINA%%%
% \newpage

\begin{center}
\Huge {\bfseries Rito de la Comuni\'on}
\end{center}

\Large {\bfseries \textcolor{red}{ANTI\'IFONA DE COMUNI\'ON \hspace{2cm} Jn 1, 12}}

\lettrine[lines=1]{\bfseries \textcolor{red}{H}}{}\Large ermanos: \\
Esta ni\~na nacida de nuevo por el Bautismo, se llama y es hija de Dios. \\
Un d\'ia recibir\'a por la Confirmaci\'on la plenitud del Esp\'iritu Santo. \\
Se acercar\'a al altar del Se\~nor, participar\'a en la mesa de su sacrificio y lo invocar\'a como Padre en medio de su Iglesia. \\
Ahora nosotros, en nombre de {\bfseries \textcolor{red}{Martina Caridad}}, \\
que es ya hija por el esp\'iritu de adopci\'on que todos hemos recibido, \\
oremos juntos como Cristo nos ense\~n\'o. \newline

\large{\bfseries \textcolor{red}{Extiende las manos y, junto con el ministro, continua:}}\newline

\Large Padre nuestro, que est\'as en el cielo...\newline

%SALTO DE PAGINA%%%
\newpage

\large{\bfseries \textcolor{red}{Solo el sacerdote, con las manos extendidas, prosigue diciendo:}}\newline

\noindent
\Large {\bfseries \textcolor{red}{V/.}} \hspace{0.5cm} L\'ibramos Se\~nor de todos lo males,\\
. \hspace{1.5cm} y conc\'edenos la paz en nuestros d\'ias,\\
. \hspace{1.5cm} para que, ayudados por tu misericordia,\\
. \hspace{1.5cm} vivamos siempre libres de pecado\\
. \hspace{1.5cm} y protegidos de toda perturbaci\'on,\\
. \hspace{1.5cm} mientras esperamos la gloriosa venida\\
. \hspace{1.5cm} de nuestro Salvador Jesucristo.\newline

\noindent
\Large {\bfseries \textcolor{red}{R/.}} \hspace{0.5cm} Tuyo es el reino, tuyo es el poder y \\
. \hspace{1.5cm} la gloria, por siempre, Se\~nor. \newline

%SALTO DE PAGINA%%%
% \newpage

\large{\bfseries \textcolor{red}{Mientras el sacerdote tama el pan consagrado, lo parte sobre la patena y pone una part\'icula dentro del caliz.}}\newline

\Large Cordero de Dios...\newline

\large{\bfseries \textcolor{red}{El sacerdote toma el pan consagrado y, sosteni\'endolo un poco elevado, dice co voz clara:}}\newline

\Large \'Este es el Cordero de Dios,\\
que quita el pecado del mundo.\\
Dichosos los invitados a la cena del Se\~nor. \newline

%SALTO DE PAGINA%%%
% \newpage

\large{\bfseries \textcolor{red}{Y, juntamente con el pueblo, a\~nade:}}\newline

\Large Se\~nor, no soy digno\\
de que entres en mi casa,
pero una palabra tuya\\
bastar\'a para sanarme.\newline

\Large {\bfseries \textcolor{red}{ORACI\'ON DESPU\'ES DE LA COMUNI\'ON}}

\lettrine[lines=1]{\bfseries \textcolor{red}{O}}{}\Large h, Dios, \\
que coronaste a la bienaventurada {\bfseries \textcolor{red}{Martina}} entre los santos \\
con el doble triunfo de la virginidad y del martirio, \\
conc\'edenos en virtud de este sacramento, \\
vencer con fortaleza toda maldad \\
y alcanzar la gloria del cielo. \\
Por Jesucristo, nuestro Se\~nor.


%SALTO DE PAGINA%%%
% \newpage

\begin{center}
\Huge {\bfseries Rito de Concluci\'on}
\end{center}

\Large {\bfseries \textcolor{red}{BENDICI\'ON SOLEMNE AL FINAL DE LA MISA}} \newline

\large{\bfseries \textcolor{red}{Seguidamente el Sacerdote bendice a la madre que tiene en sus brazos a la ni\~no(a), al padre y a todos los presentes, diciendo:}}\newline

\lettrine[lines=1]{\bfseries \textcolor{red}{E}}{}\Large l Se\~nor Todopoderoso, por su Hijo, \\
nacido de Mar\'ia la Virgen, bendiga a {\bfseries \textcolor{red}{Linda Isabel}} \\
y alegre su coraz\'on con la esperanza de la vida eterna, \\
alumbrada hoy en {\bfseries \textcolor{red}{Martina Caridad}}, \\
para que del mismo modo que le agradece el fruto de sus entra\~nas, \\ persevere con \'el en constante acci\'on de gracias. \\
Por Jesucristo nuestro Se\~nor. \newline

\noindent
\Large {\bfseries \textcolor{red}{R/.}} \hspace{0.5cm} Am\'en.\newline

%% SALTO DE PAGINA
\newpage

\noindent
\Large {\bfseries \textcolor{red}{V/.}} \hspace{0.5cm} El Se\~nor todopoderoso, \\
. \hspace{1.5cm} dispensador de la vida temporal y la eterna, \\
. \hspace{1.5cm} bendiga a {\bfseries \textcolor{red}{Samuel Jos\'e}}, \\
. \hspace{1.5cm} que junto con {\bfseries \textcolor{red}{Linda Isabel}}, \\
. \hspace{1.5cm} sean los primeros que, de palabra y obra, \\
. \hspace{1.5cm} den testimonio de la fe ante sus hijos, \\
. \hspace{1.5cm} en Jesucristo nuestro Se\~nor. \newline

\noindent
\Large {\bfseries \textcolor{red}{R/.}} \hspace{0.5cm} Am\'en.\newline

\noindent
\Large {\bfseries \textcolor{red}{V/.}} \hspace{0.5cm} El Se\~nor Todopoderoso, \\
. \hspace{1.5cm} que nos ha hecho renacer a la vida eterna \\
. \hspace{1.5cm} por el agua y el Esp\'iritu Santo, \\
. \hspace{1.5cm} bendiga a \'estos fieles, \\
. \hspace{1.5cm} para que siempre y en todo lugar sean miembros vivos de su pueblos; \\
. \hspace{1.5cm} y conceda la abundancia de su paz a todos los aqu\'i presentes, \\
. \hspace{1.5cm} en Jesucristo nuestro Se\~nor. \newline

\noindent
\Large {\bfseries \textcolor{red}{R/.}} \hspace{0.5cm} Am\'en.

\lettrine[lines=1]{\bfseries \textcolor{red}{Y}}{} \Large la bendici\'on de Dios todopoderoso, \\
Padre, Hijo \Huge{\textcolor{red}{\ding{64}}} \Large y Esp\'iritu Santo, \\
descienda sobre vosotros y os acompa\~ne siempre. \newline

\noindent
\Large {\bfseries \textcolor{red}{R/.}} \hspace{0.5cm} Am\'en.\newline

%SALTO DE PAGINA%%%
\newpage

\large{\bfseries \textcolor{red}{Luego el di\'scono, o el mismo sacerdote, con las manos juntas, despide al pueblo:}}

\lettrine[lines=1]{\bfseries \textcolor{red}{L}}{}\Large a alegr\'ia del Se\~nor sea nuestra fuerza. \\
Pod\'eis ir en paz. \newline

\noindent
\Large {\bfseries \textcolor{red}{R/.}} \hspace{0.5cm} Demos gracias a Dios. \newline

\large{\bfseries \textcolor{red}{Despu\'es el sacerdote besa con veneraci\'on el altar, como al comienzo, y, hecha la debida reverencia con los ministros, se retira a la sacrist\'ia.}}\newline

% Termina el documento %%%%%%%%%%%%%%%%%%%%%%%%%%%%% F I N %%%%%%%%%%%%%%%%%%%%%%%%%%%%%%%%%%%%%%%%%%%%%%%%%%%%%%
\end{document}
